\part{Stand}

\chapter{Stand}

\section{Der Leitfall, ein letztes Mal}

\begin{merksatz}
\boxtitel{Was der Band aus einem Wassermolekül gemacht hat}
\begin{center}
{\small
\renewcommand{\arraystretch}{1.5}
\begin{tabular}{@{}P{1.0cm}P{7.6cm}@{}}
\toprule
\textbf{Kap.} & \textbf{Was über $\mathrm{H_2O}$ gesagt wurde}\\
\midrule
1 & vier Symmetrieoperationen, Gruppe $C_{2v}\cong V_4$, vier Klassen\\
2 & neun Verschiebungskoordinaten, eine $9\times9$-Darstellung; sie ist unitär\\
3 & sie zerfällt in Irreduzible (Maschke); alle sind eindimensional (Schur, abelsch)\\
4 & ihr Charakter $(9, -1, 1, 3)$ zerlegt sich als $3A_1 + A_2 + 2B_1 + 3B_2$\\
5 & die Tafel von $C_{2v}$ aus der Gruppenstruktur allein\\
6 & drei Schwingungen $2A_1 + B_2$, alle IR- und ramanaktiv; $r_1\pm r_2$ durch Projektion\\
7 & $\mathrm{HDO}$: $C_s$, drei $A'$; die Streckungen sind $\operatorname{Ind}^{C_{2v}}_{C_s}\mathbf 1$\\
8 & zwei Protonen: ortho ($I = 1$, Gewicht $3$) und para ($I = 0$) --- Wasser hat Kernspin-Isomere\\
9 & alle Kombinationsbanden erlaubt ($5331$, $5235$ cm$^{-1}$): darum absorbiert Wasserdampf überall\\
10--11 & die $2p$-Orbitale des Sauerstoffs, $V_1\downarrow C_{2v} = A_1 + B_1 + B_2$; Spin der Protonen in $SU(2)$\\
12 & ortho/para als Schur--Weyl für $n = 2$: $S^{(2)}\otimes V_1\oplus S^{(1,1)}\otimes V_0$\\
\bottomrule
\end{tabular}}
\end{center}

\begin{center}
\begin{minipage}{0.9\textwidth}
\textbf{Ein Molekül aus drei Atomen hat den Band getragen. Alles, was über es gesagt wurde, folgt aus vier Zahlen in einer Tafel --- und die Tafel folgt aus der Gruppe. Das war die Behauptung des Vorworts, und sie ist eingelöst.}
\end{minipage}
\end{center}

\begin{center}
\begin{minipage}{0.88\textwidth}
\kursiv{\textbf{Was nicht gesagt wurde, weil die Symmetrie es nicht weiß}: die Frequenzen $3657$, $1595$, $3756$ cm$^{-1}$ (Kraftkonstanten); den Bindungswinkel $104{,}5^{\circ}$ (Elektronenstruktur); die Intensitäten (Dipolableitungen); die Wasserstoffbrücken, die aus dem Molekül die Flüssigkeit machen. Die Darstellungstheorie sagt, was null ist, was gleich ist und was sich mischen darf. Den Rest sagt die Dynamik. Ein Lehrbuch der Quantenchemie beginnt, wo dieser Band aufhört --- aber es benutzt auf jeder Seite, was hier steht. Und ortho- und para-Wasser, Zeile 8, wurden 2018 erstmals getrennt (Kilaj et al., Basel) und reagieren verschieden schnell: eine Darstellung von $S_2$, die man im Labor sortieren kann.}
\end{minipage}
\end{center}
\end{merksatz}

\section{Vier Muster}

\begin{merksatz}
\boxtitel{Muster 1 --- Entartung ist Darstellung}
Wo ein Operator mit einer Gruppe vertauscht, sind seine Eigenräume Darstellungen der Gruppe, und wenn die Gruppe mehrdimensionale irreduzible Darstellungen hat, sind Eigenwerte entartet --- nicht zufällig, sondern erzwungen. Das war die dritte Überraschung ($1, 3, 5, 7$), das $E$ des Ammoniaks, das $T_2$ des Methans, der Dirac-Punkt im Graphen, die Kramers-Dubletts, das Oktett.

\begin{center}
\begin{minipage}{0.88\textwidth}
\kursiv{\textbf{Und umgekehrt}: Mehr Entartung als die bekannte Gruppe erklärt, heißt mehr Gruppe (Wasserstoff: $SO(4)$); weniger Entartung als die Gruppe verlangt, heißt die Gruppe ist gebrochen ($\mathrm{CH_3D}$, Kristallfeld, Zeeman, Jahn--Teller). Entartungen zu zählen ist der direkteste Weg, eine Symmetrie zu finden oder zu widerlegen --- Bethe hat 1929 so die Kristallfelder gelesen, Gell-Mann 1961 so die Hadronen sortiert.}
\end{minipage}
\end{center}
\end{merksatz}

\begin{merksatz}
\boxtitel{Muster 2 --- Mitteln}
Fast jeder Beweis des Bandes hat denselben Schritt: Man nimmt ein Objekt, das nicht symmetrisch ist (ein Skalarprodukt, eine lineare Abbildung, eine Funktion), und mittelt es über die Gruppe. Das Mittel ist symmetrisch, und Schurs Lemma sagt, wie es aussieht. Satz 2.2 (Unitarität), 3.2 (Maschke), 4.3 (Orthogonalität), 6.2 (Projektion), 10.2 (kompakte Gruppen) --- alle sind Mittelung.

\begin{center}
\begin{minipage}{0.88\textwidth}
\kursiv{\textbf{Das Mittel ist das Haar-Maß}, und wo es keines gibt --- nichtkompakte Gruppen, Charakteristik $p$ mit $p\mid|G|$ ---, bricht die Theorie an genau den Stellen, an denen sie es benutzt hat: keine Unitarität, kein Maschke, kein eindeutiger Charakter, unendlichdimensionale Irreduzible. Die Bruchfälle der Kapitel 3, 4, 8 und 10 sind alle derselbe Bruchfall: Man kann nicht mitteln.}
\end{minipage}
\end{center}
\end{merksatz}

\begin{merksatz}
\boxtitel{Muster 3 --- Die Tafel reicht}
Nirgends im Band nach Kapitel 2 musste eine Matrix größer als $3\times3$ hingeschrieben werden. Charaktere sind Spuren, Spuren sind Klassenfunktionen, und Klassenfunktionen auf $D_{6h}$ sind zwölf Zahlen. Zerlegung, Auswahlregeln, Produkte, Einschränkung, Induktion, Obertöne, Indikatoren: alles Skalarprodukte von Zeilen einer Tafel.

\begin{center}
\begin{minipage}{0.88\textwidth}
\kursiv{\textbf{Was die Tafel nicht enthält}: die Basisvektoren (Clebsch--Gordan-Koeffizienten, Projektionen --- Kapitel 6 und 11 brauchten dafür die Darstellungen selbst), die Gruppe ($D_4$ gegen $Q_8$; aber Tafel plus Quadratabbildung reicht), und für $S_n$ über $\mathbb F_p$ nicht einmal die Dimensionen. Die Tafel ist die Darstellungstheorie, soweit sie kommutativ ist --- und das ist, für endliche Gruppen über $\mathbb C$, erstaunlich weit.}
\end{minipage}
\end{center}
\end{merksatz}

\begin{merksatz}
\boxtitel{Muster 4 --- Gebrochen ist informativer als ganz}
Die Stellen, an denen die Symmetrie nicht ganz gilt, haben jedes Mal die Physik getragen: die verbotene Bande bei $255$ nm (Herzberg--Teller), die Fermi-Resonanz, das Kristallfeld, der Zeeman-Effekt, Jahn--Teller, die Quarkmassen, die Quasikristalle. Eine exakte Symmetrie liefert Entartungen und Nullen; eine leicht gebrochene liefert Aufspaltungen, deren Muster die ursprüngliche Gruppe verrät und deren Größe die Störung misst.

\begin{center}
\begin{minipage}{0.88\textwidth}
\kursiv{\textbf{Das ist die Methode der Störungstheorie}, und die Darstellungstheorie ist ihr Ordnungsprinzip: Die Störung transformiert sich nach irgendeiner Darstellung ($E_g$ für die tetragonale Verzerrung, $V_4$ für das Oktaederfeld, die $Y$-Komponente der $\mathbf 8$ für die $s$-Masse), und welche Niveaus sie wie koppelt, steht in den Clebsch--Gordan-Reihen. Wigner und Eckart haben das 1927--30 als Satz formuliert: Alle Matrixelemente eines Tensoroperators zwischen Zuständen fester Symmetrie sind ein einziges \enquote{reduziertes} Matrixelement mal Clebsch--Gordan-Koeffizienten. Der Band hat den Satz benutzt, ohne ihn auszusprechen --- in jeder Auswahlregel.}
\end{minipage}
\end{center}
\end{merksatz}

\section{Was der Band ausgelassen hat}

\begin{merksatz}
\boxtitel{Die Landkarte jenseits dieses Bandes}
\begin{center}
{\small
\renewcommand{\arraystretch}{1.45}
\begin{tabular}{@{}P{3.2cm}P{5.4cm}@{}}
\toprule
\textbf{Lie-Algebren systematisch} & Kapitel 11 hat $\mathfrak{su}(2)$ mit Leiteroperatoren behandelt und $SU(3)$ nur benutzt. Die Klassifikation der halbeinfachen Lie-Algebren (Killing, Cartan: $A_n, B_n, C_n, D_n, G_2, F_4, E_6, E_7, E_8$), Wurzelsysteme, Weyls Charakterformel, Dynkin-Diagramme --- das ist ein eigener Band.\\
\textbf{Unendlich\-dimensionale Darstellungen} & Die Poincaré-Gruppe und Wigners Klassifikation der Teilchen (Masse, Spin), Harish-Chandras Theorie, Kac--Moody-Algebren und die Virasoro-Algebra der konformen Feldtheorie: nur als Bruchfall erwähnt.\\
\textbf{Modulare Darstellungstheorie} & Brauer-Charaktere, Blöcke, Defektgruppen; die Theorie, in der die Bruchfälle der Kapitel 3, 4 und 8 leben. Mit $S_n$ über $\mathbb F_p$ als offenem Zentrum.\\
\textbf{Darstellungen über $\mathbb Q$ und $\mathbb R$} & Schur-Index, Rationalitätsfragen; der Band hat nur den Indikator (Kapitel 9) gegeben.\\
\bottomrule
\end{tabular}}
\end{center}
\begin{center}
{\small
\renewcommand{\arraystretch}{1.45}
\begin{tabular}{@{}P{3.2cm}P{5.4cm}@{}}
\toprule
\textbf{Kategorien} & Tensorkategorien, Fusionskategorien, Tannaka-Dualität (die Gruppe aus ihrer Darstellungskategorie zurückgewinnen), Kategorifizierung; Band 23 hat die Sprache, dieser Band hat sie nicht benutzt.\\
\textbf{Algorithmen} & Dixon--Schneider, Meataxe, Unger; die Programme GAP und Magma; das Programm dieses Bandes (\texttt{rep.py}) rechnet Tafeln aus Matrizen, nicht aus der Gruppenstruktur.\\
\textbf{Raumgruppen vollständig} & Die $230$ wurden gezählt, nicht hergeleitet; die Darstellungstheorie der Raumgruppen (kleine Gruppen, projektive Darstellungen bei nichtsymmorphen Gruppen) nur angedeutet.\\
\textbf{Langlands} & Das Programm, das Darstellungen von Galois-Gruppen mit automorphen Darstellungen verbindet und die halbe Zahlentheorie seit 1970 organisiert: nicht angerührt, nur in Kapitel 7 als Fernziel der Induktion genannt.\\
\bottomrule
\end{tabular}}
\end{center}
\end{merksatz}

\section{Was offen ist}

\begin{merksatz}
\boxtitel{Stand bei Redaktionsschluss, August 2026}
\begin{center}
{\small
\renewcommand{\arraystretch}{1.45}
\begin{tabular}{@{}P{3.2cm}P{5.4cm}@{}}
\toprule
\textbf{Gelöst in den letzten Jahren} & McKay-Vermutung (Cabanes--Späth 2024, nach Reduktion auf einfache Gruppen durch Isaacs--Malle--Navarro 2007); umbrales Moonshine (Duncan--Griffin--Ono 2015); geometrisches Langlands über Funktionenkörpern (Gaitsgory, Raskin et al., angekündigt 2024, fünf Arbeiten, rund $1000$ Seiten); die Charaktertafel-Frage nach Elementordnungen für $p$-Gruppen (Navarro et al.).\\
\textbf{Offen: Zählprobleme} & Alperins Gewichtsvermutung (1987), Brauers $k(B)$-Vermutung (1956, bewiesen für $p$-auflösbare Gruppen 2004), Broués Vermutung über abelsche Defektgruppen (1990): alle sagen, dass die Zahl der Charaktere in einem Block \enquote{lokal} bestimmt ist. Reduziert auf einfache Gruppen, dort in Arbeit.\\
\textbf{Offen: Kronecker} & Kombinatorische Formel für $g_{\lambda\mu\nu}$ (seit 1938); Plethysmus-Koeffizienten; Saxl-Vermutung (gibt es ein $S^{\lambda}$, dessen Quadrat alle Irreduziblen enthält? für Treppenpartitionen vermutet 2012).\\
\bottomrule
\end{tabular}}
\end{center}
\begin{center}
{\small
\renewcommand{\arraystretch}{1.45}
\begin{tabular}{@{}P{3.2cm}P{5.4cm}@{}}
\toprule
\textbf{Offen: modular} & Dimensionen der irreduziblen $\mathbb F_p S_n$-Moduln; nach dem Fall der James-Vermutung keine Ersatzvermutung in Sicht. Lusztigs Charakterformel für $GL_n(\mathbb F_p)$ gilt für $p\gg n$, aber die Schranke ist unbekannt (Williamson: mindestens exponentiell).\\
\textbf{Offen: Analysis} & Das unitäre Dual nichtkompakter Gruppen: für $E_8$ 2007 berechnet (Atlas of Lie Groups, $453\,060$ Darstellungen), allgemein ohne Klassifikation; Arthurs Vermutungen; die Plancherel-Formel für $p$-adische Gruppen in Spezialfällen.\\
\textbf{Offen: Algorithmen} & Schnelle Fourier-Transformation in $O(|G|\log^{c}|G|)$ für alle endlichen Gruppen; Graphisomorphie als verstecktes Untergruppenproblem (Babai 2015: quasipolynomial klassisch, Quantenvorteil ungeklärt); Kronecker-Positivität in $\#P$.\\
\textbf{Offen: Physik} & Ob die Jahn--Teller-Aussage einen fallfreien Beweis hat; ob $SU(3)_{\mathrm{Flavour}}$-Brechung aus der QCD quantitativ folgt (Gitter: ja, zu $5\,\%$); die Klassifikation topologischer Phasen mit Raumgruppensymmetrie jenseits freier Elektronen.\\
\bottomrule
\end{tabular}}
\end{center}

\begin{center}
\begin{minipage}{0.88\textwidth}
\kursiv{\textbf{Die größte Bewegung der letzten Jahre} ist die Begegnung zwischen Darstellungstheorie und maschinellem Lernen: äquivariante Netze (Kapitel 10) als Konstruktionsprinzip, und umgekehrt Netze, die Vermutungen über Kazhdan--Lusztig-Polynome gefunden haben (Davies et al., \textit{Nature} 2021, mit Williamson) --- eine neue Formel für die Polynome über die \enquote{Hyperwürfel} eines Bruhat-Intervalls, von Menschen bewiesen, von einem Netz gefunden. Ob daraus eine Methode wird, ist im August 2026 offen; dass sie funktioniert hat, nicht.}
\end{minipage}
\end{center}
\end{merksatz}

\section{Sechs Sätze, die bleiben}

\begin{merksatz}
\boxtitel{Die Bilanz}
\begin{center}
{\small
\renewcommand{\arraystretch}{1.6}
\begin{tabular}{@{}P{3.2cm}P{5.4cm}@{}}
\toprule
\textbf{Maschke (3.2)} & Jede Darstellung einer endlichen Gruppe über $\mathbb C$ ist eine direkte Summe irreduzibler.\\
\textbf{Schur (3.3)} & Zwischen Irreduziblen gibt es nur Null und Isomorphismen; Entartung ist Darstellung.\\
\textbf{Orthogonalität (4.3, 5.2)} & Die irreduziblen Charaktere sind orthonormal; die Tafel ist unitär; $m_i = \langle\chi, \chi_i\rangle$.\\
\textbf{Zählung (4.5)} & So viele Irreduzible wie Klassen; $\sum d_i^{2} = |G|$.\\
\textbf{Reziprozität (7.5)} & $\langle\operatorname{Ind}\sigma, \rho\rangle = \langle\sigma, \operatorname{Res}\rho\rangle$.\\
\textbf{Peter--Weyl (12.3)} & Für kompakte Gruppen bleibt alles, mit Integral statt Summe; die Matrixelemente sind vollständig.\\
\bottomrule
\end{tabular}}
\end{center}

\begin{center}
\begin{minipage}{0.9\textwidth}
\textbf{Sechs Sätze, alle vor 1927 bewiesen, alle mit Mittelung bewiesen. Was der Band aus ihnen gemacht hat --- Schwingungsspektren, Orbitale, Spin, Kristallfelder, Hadronen ---, ist Anwendung, nicht Erweiterung. Die Theorie ist seit hundert Jahren fertig; die Natur hört nicht auf, sie zu benutzen.}
\end{minipage}
\end{center}
\end{merksatz}

\begin{merksatz}
\boxtitel{Ein Schlusswort}
\begin{center}
\begin{minipage}{0.9\textwidth}
\kursiv{Frobenius hat 1896 Charaktere erfunden, um eine Determinante zu faktorisieren, die Dedekind ihm in einem Brief geschickt hatte. Dreißig Jahre später erklärten dieselben Charaktere, warum Atome Spektren haben, die sie haben; vierzig Jahre später, warum Benzol stabil ist; siebzig Jahre später, warum es ein Teilchen namens $\Omega^{-}$ geben muss. Nichts davon war beabsichtigt. Die Darstellungstheorie ist das deutlichste Beispiel in der Mathematik dafür, dass eine Antwort auf eine Frage, die niemand gestellt hat, die Fragen von drei Wissenschaften beantwortet.}

\medskip
\kursiv{Was sie liefert, ist bescheiden formuliert: welche Zahlen null sind, welche gleich sind und was sich mischen darf. Aber das ist genau die Information, die man braucht, bevor man rechnet --- und oft die einzige, die man braucht, um zu verstehen, was man gemessen hat. Ein Chemiker, der die Tafel von $C_{2v}$ kennt, weiß über das Wasser Dinge, die ohne sie nur ein Computer herausfände. Ein Physiker, der weiß, dass $\mathrm{Sym}^{2}V_{1}$ kein $V_1$ enthält, weiß, dass das Higgs-Teilchen nicht Spin $1$ hat, bevor irgendjemand eine Winkelverteilung gemessen hat.}

\medskip
\kursiv{\textbf{Dieser Band hat versucht, nichts für selbstverständlich zu halten}: nicht, dass Spuren etwas bedeuten; nicht, dass Entartungen Gründe haben; nicht, dass eine Tafel aus Zahlen ein Molekül beschreiben kann. Wenn er gelungen ist, sieht der Leser in der Tafel von $D_{6h}$ nicht zwölf mal zwölf Zahlen, sondern das Benzol --- und in der Tafel von $C_{2v}$ das Wasser, mit allem, was es kann.}
\end{minipage}
\end{center}
\end{merksatz}
